% ! TeX root = ./main_ieee.tex
\section{Evaluation of Placement Policies}\label{sec:amr:eval}

\newcommand{\baseline}{\texttt{Baseline}} \newcommand{\cpla}{\texttt{CPL0}}
\newcommand{\cplb}{\texttt{CPL25}} \newcommand{\cplc}{\texttt{CPL50}}
\newcommand{\cpld}{\texttt{CPL75}} \newcommand{\cple}{\texttt{CPL100}}

\tableruns

In this section, we evaluate CPLX under real AMR workloads and synthetic microbenchmarks.
While {\cple} (pure LPT) and the baseline policy represent known quantities, our key question is
whether intermediate values of $X$ can provide meaningful control over the locality-balance
tradeoff.
For CPLX to be successful, it must achieve three goals: intermediate values should demonstrate
measurable control over both communication patterns and runtime performance, achieve minimum
runtime if the tradeoff is worth balancing, and maintain low runtime overhead even in setups with
rapid refinement.
We evaluate these aspects by comparing different values of $X$ against {\cple} as our reference
point.

We primarily evaluate CPLX using the Sedov Blast Wave 3D problem in
Phoebus~\cite{Barke2024:Phoebus}, a hydrodynamics code built on top of the
Parthenon~\cite{Grete2023:Parthenon} framework.
While we also studied placement in other codes, such as a galaxy cooling setup in
AthenaPK~\cite{Grete2023:Parthenon}, results were directionally similar: codes with high compute
variability benefit more from better placement, and vice-versa.
We describe our hardware setup in \S\ref{sec:amr:eval:setup}, present Sedov results in
\S\ref{sec:amr:eval:sedov}, and use microbenchmarks to evaluate CPLX's load balance and locality
properties under different synthetic regimes in \S\ref{sec:amr:eval:micro}.


\subsection{Experimental Setup}\label{sec:amr:eval:setup}

% We use four setups of the Sedov Blast Wave problem, as shown in Table 1. We configure the initial
% mesh size to initially have one block per rank. The number of blocks quickly increases due to
% refinement, settling at $\sim$2 blocks per rank towards the end of the simulation. The initial and
% final block counts, and the total and the load-balancing timesteps are also shown in Table 1. 
We evaluate four configurations of the Sedov Blast Wave problem, as detailed in Table 1.
Each run begins with one block per rank, increasing as refinement proceeds, and finally settling at
\mytilde{}two blocks per rank.
The table reports initial and final block counts, total steps, and load-balancing invocations.

% All experiments were run on the tuned research cluster described in \S\ref{sec:amr:tele}, with 16
% ranks per node, at scales ranging from 512--4096 ranks. All telemetry was collected using the
% custom profiling tool described in \S\ref{sec:amr:chal}. We compare the baseline policy with 5
% parameters for CPLX (0--100 in increments of 25). The results reported here are gains from
% placement against the tuned baseline. We do not report runtime gains from tuning because of the
% variable impact of different parameters across codes and scales, and the difficulty of reproducing
% issues with fail-slow hardware.

All experiments ran on the tuned Emulab-backed cluster described in \S\ref{sec:amr:tele}, with 16 ranks
per node, at scales from 512--4096 ranks.
Telemetry was collected using the custom profiling stack described in \S\ref{sec:amr:chal}.
We compare the tuned baseline against CPLX with X $\in \{0, 25, 50, 75, 100\}$, and only report
runtime gains from placement.
Runtime gains from tuning are omitted because of variable impact across different codes and scales.

\parahead{Hardware Checks}
To ensure experimental integrity, all reserved nodes were scanned for hardware issues---thermal
throttling, memory errors etc.---before every individual run in each parameter suite using the
process described in \S\ref{sec:amr:tele}.
All reported results are from runs where no participating node showed issues either before or after
execution.

\parahead{Software Stack}
All experiments used MVAPICH2~v2.3.7~\cite{Panda2021:MVAPICH}, compiled against a tuned version of
PSM~\cite{2020:IntelPerformanceScaled}, the userspace library for our QLogic fabric.

\figAmrEvalRuntime

\subsection{Sedov Blast Wave Results}\label{sec:amr:eval:sedov}

We evaluate five placement policies on the Sedov Blast Wave problem across four scales (512--4096
ranks), comparing baseline performance against CPLX, as $X$ is varied from 0--100.
% : baseline (CPL0), pure LPT (CPL100), and three CPLX hybrids with intermediate values of the
% locality parameter $X$. All runs were executed on a tuned cluster with verified measurement
% stability.
The analysis focuses on overall runtime behavior (\cref{fig:amr:eval:runtime:wide}), the
synchronization---communication tradeoff (\cref{fig:amr:eval:tradeoff:wide}), and P2P communication
patterns (\cref{fig:amr:eval:comm:wide}).

\finding{1}{Baseline: synchronization reaches 50\% of total runtime at scale}

\noindent\Cref{fig:amr:eval:runtime:wide} shows total runtime across all policies,
decomposed into computation, communication, synchronization, and rebalancing
phases.
With baseline placement, computation and synchronization dominate the profile, jointly accounting
for over 90\% of execution time at all scales.
Communication and rebalancing remain minor components, contributing approximately 7\% and 3\%
respectively.
Synchronization overhead grows sharply with scale---from 35\% at 512 ranks to 50\% at 4096 ranks,
despite the tuned environment.
These results demonstrate the challenge of managing dynamic behavior in bulk-synchronous parallel
applications.
Unless mitigated, runtime is dictated by stragglers, the likelihood of which only goes up with
scale.

\finding{2}{CPLX achieves up to 21.6\% overall runtime reduction, with impact increasing with scale}

\noindent\Cref{fig:amr:eval:runtime:wide} also shows that all CPLX configurations outperform
the baseline, with runtime improvements that grow more pronounced at larger scales.
At 4096 ranks, the best-performing variant (\cple) reduces total runtime by up to \textbf{21.6\%}
relative to baseline.
Even at 512 ranks, CPLX achieves a 15.3\% reduction, and all values of $X$ tested yield
improvements above 12\%.

Runtime follows a clear U-shaped curve as a function of $X$: starting from \cpla{}
(locality-preserving), decreasing to a minimum with intermediate values, and rising again toward
\cple{} (pure LPT).
This shape reflects the ability of CPLX to control the load-locality tradeoff in a manner that
yields meaningful overall runtime gains.

% This shape reflects the realized runtime effect of competing effects of improved load balance
% (which reduces synchronization) and degraded locality (which increases communication), with
% intermediate values of $X$ achieving the best compromise. 
Crucially, compute time remains flat across all policies, confirming that the total amount of work
is invariant to placement.
The observed improvements arise entirely from reductions in synchronization overhead, which
dominate non-compute runtime in the baseline configuration.
Measured as reduction in non-compute time, the impact of CPLX becomes even more pronounced---up to
36\% reduction at 4096 ranks.

\figAmrEvalTradeoff

\finding{3}{CPLX enables consistent control over the load-locality tradeoff, and tradeoff impact increases with scale}

\noindent\Cref{fig:amr:eval:tradeoff:wide} isolates the core tradeoff by showing P2P communication
and synchronization times, normalized against the baseline, for all policy configurations at 512 and 4096 ranks.
Intermediate scales are omitted for compactness, but we find them to demonstrate similar trends.
Increasing $X$ reduces synchronization time, as expected from better load balance, while increasing
communication time due to loss of locality.

We find that modest values of $X$ (25--50) capture almost all of LPT load-balancing benefits,
without incurring the commensurate locality cost.
While LPT incurs a relative increase of 55\% in locality cost, its impact on the overall runtime
remains modest as even at 4096 ranks, communication is only 13\% of the overall runtime.
The increasing spread between intermediate $X$ values and LPT with scale suggests that the runtime
impact of balancing the two increases with scale.
Given the strong performance of LPT vs commercial ILP solvers (\S\ref{sec:amr:des:lpt}), this is likely
close to the practically achievable optimum under our model of AMR placement.

At 512 ranks, some CPLX configurations even reduce communication time relative to baseline---this
is explained by the compounding effects captured by boundary communication time.
High variance in the preceding compute kernels results in more time spent in \mpiw{} by some ranks,
and improving load-balance has a positive downstream impact that overrides the locality cost.
This is an illustration of higher-order placement effects that we discard for computational
tractability.

\figAmrEvalComm

\finding{4}{P2P message volume statistics confirm locality degradation with increasing $X$}

\noindent\cref{fig:amr:eval:comm:wide} shows the breakdown of local (intra-node shared memory) and
remote (inter-node MPI) P2P messages, normalized to the total baseline message volume,
for \cplx{} variants at 512 and 4096 ranks.
As $X$ increases, remote message counts rise while local messages fall, reflecting the fact that
CPLX progressively expands LPT coverage at the expense of locality.
This tradeoff is enacted mechanically by CPLX, independent of the underlying application or
mesh---only the runtime impact of this shift depends on workload characteristics.
We explore workload variations further in \S\ref{sec:amr:eval:micro} using synthetic microbenchmarks.

A subtler effect visible in \cref{fig:amr:eval:comm:wide} is the growth in total MPI-visible message
volume with increasing $X$.
This occurs because intra-rank communication---handled via \texttt{memcpy} when blocks are
co-located---is replaced by MPI messages when placement breaks locality.
% When placement breaks locality, these transfers become inter-rank messages, now visible in
% MPI-level statistics.
This shift explains part of the communication time increase observed in
\cref{fig:amr:eval:tradeoff:wide}.

Finally, even the baseline placement routes a majority of messages across nodes: at 4096 ranks,
\textbf{64\%} of messages are already remote.
This is an intrinsic property of space-filling curves---dimensionality reduction to 1-D preserves
only partial spatial locality.
\cplx{} does not introduce a new communication cost class, but extends an existing placement pattern and exposes it as a tunable parameter.
% It extends an existing tradeoff and exposes it as a tunable axis for balancing locality against
% synchronization overhead.

\subsection{Microbenchmarks}\label{sec:amr:eval:micro}

To complement the application-level Sedov results (\S\ref{sec:amr:eval:sedov}) and enable controlled
evaluation across a broader range of conditions, we developed two microbenchmarks.
\commbench{} isolates P2P communication under varying placements, while \scalebench{} evaluates the effectiveness and
computational cost of placement policies under synthetic compute imbalance.

\figAmrMicroLoc

\parahead{Commbench: Simulating Boundary Communication Patterns}

% \commbench{} is a synthetic microbenchmark designed to isolate boundary communication patterns and
% evaluate the impact of placement locality on end-to-end round latency. It constructs octree-based
% AMR meshes with realistic refinement, and derives P2P communication patterns from geometric
% neighbor relationships (face, edge, vertex) tracked across refinement levels. \commbench{} is
% inspired by the Halo3D-26 benchmark, but instead of using fixed-grid meshes with fixed placements,
% it simulates a complete AMR placement pipeline and supports custom placement policies as drop-in
% modules.

\commbench{} is a synthetic microbenchmark that isolates boundary communication and
evaluates how placement locality affects end-to-end round latency.
It constructs octree-based AMR meshes with realistic refinement, deriving P2P patterns from
geometric neighbor relationships (face, edge, vertex) across refinement levels.
While inspired by the Halo3D-26 fixed-grid benchmark~\cite{Rodri2011:SST}, \commbench{} simulates
a full AMR placement pipeline and accepts custom placement policies as drop-in modules.

Each communication round emulates a boundary exchange, with blocks exchanging messages based on
neighbor type.
Message sizes are realistic: face-neighbor exchanges are proportionally larger than edge or vertex
ones.
Meshes are refined to yield 1--2 blocks per rank, and each round includes all neighbor types.
Timing is captured using MPI barriers before and after each round.
Results are averaged over 100 rounds and 10 random meshes per policy.
We discard cold-start rounds and those with latency above 10\,ms, which we found to reflect
fabric-level recovery behavior unrelated to placement.

% Each communication round emulates a boundary exchange, with blocks exchanging messages based on
% their geometric neighbor types. Message sizes are modeled realistically: face-neighbor exchanges
% are larger than edge or vertex ones. Meshes are refined to yield 1--2 blocks per rank, and each
% round includes all neighbor types. Timing is captured using MPI barriers before and after each
% round. Results are averaged over 100 rounds and 10 random meshes per policy. We discard cold-start
% rounds as well as those with latency above 10\,ms, which we found to reflect fabric-level recovery
% behavior that obscures placement-driven trends.

\Cref{fig:amr:eval:loc} shows average round latency across placement policies as
locality decreases from \cpla{} to \cple.
While higher locality yields lower latency at small scales (512 ranks and below), a surprising
U-shaped trend emerges at larger scales: intermediate values of $X$ outperform both extremes.
This inversion is driven by face-neighbor communication, which increases as meshes get denser.
Locality-preserving policies cluster high-traffic neighbors unevenly, increasing per-rank load.
Intermediate \cplx{} settings diffuse this clustering without fully randomizing placement, leading
to more uniform traffic patterns.
Though the latency differences are modest ($\pm$0.5\,ms), the reversal of expected trends
underscores how observed system behavior can defy theoretical intuitions.
\commbench{} provides a practical mechanism for empirically selecting $X$ in different environments.

\figAmrMicroScalA
\figAmrMicroScalB

\parahead{Scalebench: Evaluating Imbalance and Placement Overhead}
\scalebench{} evaluates different policies at scales from 512 to 128K ranks, with 1--2 blocks per rank to enable flexible load-balancing.
Block costs are drawn from three representative distributions---exponential, Gaussian, and
power-law---with variability bounds chosen to create meaningful balancing opportunities while
remaining within realistic AMR ranges.

\Cref{fig:amr:eval:scal_a} plots the normalized makespan across placement policies (lower is better).
We see that \cple{} (LPT) achieves the lowest makespan across distributions, but \cpla{} and
\cplb{} capture the bulk of the benefits with much higher locality retention.
\Cref{fig:amr:eval:scal_b} reports placement computation overhead as a function of scale.
Compute costs increase with scale, but remain tractable (\mytilde{}10\,ms) up to 16K ranks, rising
to 100\,ms at 128K ranks.
At the largest scales, zonal placement architectures can be adopted to mitigate placement
overhead---dividing ranks into $k$ zones to compute placement independently and in
parallel~\cite{Zheng2011:PeriodicHierarchicalLoad}.


% Loaded data for runs: nd512.2, nd1024.3, nd2048.2, nd4096.4 Analyzing run: nd512.2
% ========================== Phase order: ['Total', 'Compute', 'Communication', 'Synchronization',
% 'Rebalancing'] Expected order: ['Total', 'Compute', 'Communication', 'Synchronization',
% 'Rebalancing'] Total : 8458s (100.00%), 7430s (100.00%), 7241s (100.00%), 7163s (100.00%), 7240s
% (100.00%), 7357s (100.00%) Compute : 4692s (55.47%), 4694s (63.18%), 4696s (64.85%), 4701s
% (65.63%), 4703s (64.96%), 4707s (63.98%) Communication : 668s (7.90% ), 637s (8.57% ), 638s (8.81%
% ), 640s (8.93% ), 680s (9.39% ), 781s (10.62%) Synchronization: 3028s (35.80%), 2028s (27.29%),
% 1834s (25.33%), 1746s (24.38%), 1777s (24.54%), 1791s (24.34%) Rebalancing : 70s (0.83% ), 71s
% (0.96% ), 73s (1.01% ), 76s (1.06% ), 80s (1.10% ), 78s (1.06% ) Analyzing run: nd1024.3
% ========================== Phase order: ['Total', 'Compute', 'Communication', 'Synchronization',
% 'Rebalancing'] Expected order: ['Total', 'Compute', 'Communication', 'Synchronization',
% 'Rebalancing'] Total : 12764s (100.00%),10770s (100.00%),10671s (100.00%),10452s (100.00%),10604s
% (100.00%),10905s (100.00%) Compute : 6446s (50.50%), 6465s (60.03%), 6448s (60.43%), 6451s
% (61.72%), 6448s (60.81%), 6508s (59.68%) Communication : 909s (7.12% ), 948s (8.80% ), 960s (9.00%
% ), 984s (9.41% ), 1050s (9.90% ), 1253s (11.49%) Synchronization: 5175s (40.54%), 3115s (28.92%),
% 3001s (28.12%), 2742s (26.23%), 2819s (26.58%), 2844s (26.08%) Rebalancing : 234s (1.83% ), 242s
% (2.25% ), 262s (2.46% ), 275s (2.63% ), 287s (2.71% ), 300s (2.75% ) Analyzing run: nd2048.2
% ========================== Phase order: ['Total', 'Compute', 'Communication', 'Synchronization',
% 'Rebalancing'] Expected order: ['Total', 'Compute', 'Communication', 'Synchronization',
% 'Rebalancing'] Total : 9719s (100.00%), 8039s (100.00%), 7739s (100.00%), 7712s (100.00%), 7737s
% (100.00%), 8049s (100.00%) Compute : 4382s (45.09%), 4396s (54.68%), 4389s (56.71%), 4389s
% (56.91%), 4389s (56.73%), 4419s (54.90%) Communication : 709s (7.29% ), 790s (9.83% ), 804s
% (10.39%), 818s (10.61%), 839s (10.84%), 1045s (12.98%) Synchronization: 4462s (45.91%), 2677s
% (33.30%), 2364s (30.55%), 2321s (30.10%), 2318s (29.96%), 2377s (29.53%) Rebalancing : 166s (1.71%
% ), 176s (2.19% ), 182s (2.35% ), 184s (2.39% ), 191s (2.47% ), 208s (2.58% ) Analyzing run:
% nd4096.4 ========================== Phase order: ['Total', 'Compute', 'Communication',
% 'Synchronization', 'Rebalancing'] Expected order: ['Total', 'Compute', 'Communication',
% 'Synchronization', 'Rebalancing'] Total : 12403s (100.00%),10311s (100.00%), 9924s (100.00%),
% 9730s (100.00%), 9748s (100.00%),10352s (100.00%) Compute : 5012s (40.41%), 5048s (48.96%), 5013s
% (50.51%), 5006s (51.45%), 5006s (51.35%), 5072s (49.00%) Communication : 871s (7.02% ), 971s
% (9.42% ), 987s (9.95% ), 1010s (10.38%), 1036s (10.63%), 1347s (13.01%) Synchronization: 6204s
% (50.02%), 3925s (38.07%), 3556s (35.83%), 3331s (34.23%), 3313s (33.99%), 3525s (34.05%)
% Rebalancing : 316s (2.55% ), 367s (3.56% ), 368s (3.71% ), 383s (3.94% ), 393s (4.03% ), 408s
% (3.94% ) -------------------------- 
